\paragraph{Observações.} Com ar passivo, (iii) e (ii) dão campos quase iguais: no f50 com 36~mm/h a zona saturada
aflora na face e resta sucção sob a crista; com 90~mm/h todo o talude tem $p_w \ge 0$ (o lençol chegou à superfície).
Com ar ativo (i) o quadro é outro: o ar comprimido (até 58~kPa no f50, o painel de $p_a$) eleva $p_w$, e em 24~h não
resta sucção nem sob a crista, embora tenha entrado menos água. No f100 com 9~mm/h a água fica numa camada de 1--2~m
junto à superfície, com a sucção inicial preservada abaixo nas variantes de ar passivo; essa camada tem a espessura de
um ou dois elementos da malha do transiente, o que pede o teste de malha.
